\documentclass[10pt]{article}
\usepackage[T1]{fontenc}
\usepackage[utf8]{inputenc}
\usepackage[french]{babel}
\usepackage{amsmath,amssymb,mathrsfs,mathpazo}
\usepackage{geometry}
\geometry{a4paper,margin=2cm}
\usepackage{graphicx,array,tabularx,tabularray,colortbl,multirow,multicol,booktabs,enumitem,pifont,ulem,fancybox,diagbox,cancel}
\usepackage{pstricks,pstricks-add,pst-plot,pst-tree,pst-node,pst-text,pst-func,pst-math,pst-bspline,pst-3dplot,pst-eucl,pst-solides3d}
\let\pscalculate\undefined
\usepackage{pst-bezier}
\usepackage{tikz,pgfplots,tkz-tab}
\usetikzlibrary{arrows,arrows.meta,patterns,calc,positioning,trees,decorations.markings,shapes.geometric,shapes.misc,shapes.arrows,intersections,decorations.pathmorphing,decorations.pathreplacing,backgrounds,shadows,babel}
\pgfplotsset{compat=1.18}
\usepgfplotslibrary{fillbetween}
\usepackage[np]{numprint}
\usepackage{listings,makecell,pas-tableur,twoopt}
\let\barre\undefined
\usepackage[pstricks]{tablvar}
\definecolor{darkgreen}{RGB}{0,100,0}
\providecommand{\red}{\color{red}}
\providecommand{\blue}{\color{blue}}
\providecommand{\green}{\color{green}}
\providecommand{\black}{\color{black}}
\providecommand{\white}{\color{white}}
\providecommand{\gray}{\color{gray}}
\providecommand{\cyan}{\color{cyan}}
\providecommand{\magenta}{\color{magenta}}
\providecommand{\yellow}{\color{yellow}}
\providecommand{\violet}{\color{violet}}
\providecommand{\orange}{\color{orange}}
\providecommand{\euro}{\texteuro}
\providecommand{\EUR}{\texteuro}
\providecommand{\ieme}{\textsuperscript{e}}
\providecommand{\degres}{\ensuremath{^\circ}}
\providecommand{\degree}{\ensuremath{^\circ}}
\providecommand{\vv}[1]{\overrightarrow{#1}}
\providecommand{\Coord}[1]{#1}
\providecommand{\decosix}{}
\providecommand{\starredbullet}{\ensuremath{\bullet}}
\pagestyle{empty}
\setlength{\parindent}{0pt}
\usepackage{adjustbox}\DeclareUnicodeCharacter{2236}{:}\DeclareUnicodeCharacter{2190}{\ensuremath{\leftarrow}}
\providecommand{\R}{\mathbb{R}}
\providecommand{\N}{\mathbb{N}}
\providecommand{\Z}{\mathbb{Z}}
\providecommand{\Q}{\mathbb{Q}}
\providecommand{\D}{\mathbb{D}}
\providecommand{\C}{\mathbb{C}}
\providecommand{\e}{{\mathrm{e}}}
\providecommand{\vect}[1]{\overrightarrow{#1}}
\providecommand{\vectt}[1]{\overrightarrow{\mathrm{#1}}}
\providecommand{\vvt}[1]{\overrightarrow{\mathrm{#1}}}
\providecommand{\np}[2][]{{#2~#1}}
\providecommand{\vv}[1]{{\overrightarrow{#1}}}
\renewcommand{\d}{{\mathrm{d}}}
\providecommand{\degres}{{^\circ}}
\providecommand{\degre}{{^\circ}}
\providecommand{\blue}{\color{blue}}
\providecommand{\red}{\color{red}}
\providecommand{\enskip}{\quad}
\providecommand{\llbracket}{\lbrack\!\lbrack}
\providecommand{\rrbracket}{\rbrack\!\rbrack}
\providecommand{\euro}{{\text{€}}}
\providecommand{\no}{{\text{n°}}}
\providecommand{\interval}[2]{{[#1;#2]}}
\providecommand{\intervFO}[2]{{[#1;#2[}}
\providecommand{\dotfill}{\ldots}

\providecommand{\ds}{\displaystyle}
\providecommand{\textding}[1]{#1}
\providecommand{\Oij}{\(\mathrm{O};\vec{\imath},\vec{\jmath}\)}
\providecommand{\Oijk}{\(\mathrm{O};\vec{\imath},\vec{\jmath},\vec{k}\)}
\providecommand{\Ouv}{\(\mathrm{O};\vec{u},\vec{v}\)}
\providecommand{\intervFF}[2]{\ensuremath{[#1;#2]}}
\renewcommand{\euro}{\ifmmode\text{\texteuro}\else\texteuro\fi}
\DeclareUnicodeCharacter{00B0}{\ifmmode^{\circ}\else\textdegree\fi}
\DeclareUnicodeCharacter{20AC}{\ifmmode\text{\texteuro}\else\texteuro\fi}
% Énoncés : Éducation nationale. Transcription et corrigés : APMEP, Denis Vergès et contributeurs. Découpage et classement : CoopMaths, Sébastien Lozano, Éric Elter et contributeurs. Adaptation de la présentation : Maths Explorer.
% CC BY-SA : https://coopmaths.fr/www/about/
% https://coopmaths.fr/alea/static/bac/2023/tex/bac_2023_03_sujet2_polynesie_1.tex
\begin{document}
\noindent\textit{APMEP / CoopMaths. CC BY-SA. Présentation adaptée par Maths Explorer.}\par\medskip
\medskip

\emph{Les parties {\rm A} et {\rm B} peuvent être traitées indépendamment}

\medskip

\textbf{Partie A}

\medskip
Chaque jour, un athlète doit sauter une  haie en fin d'entraînement. Son entraîneur estime, au vu de la saison précédente que

\begin{itemize}
\item[$\bullet~~$]si l'athlète franchit la haie un jour,  alors il la franchira dans $90\,\%$ des cas le jour suivant ;
\item[$\bullet~~$]si l'athlète ne franchit pas la haie un jour, alors dans $70\,\%$ des cas il ne la franchira pas non plus le lendemain.
\end{itemize}

On note pour tout entier naturel $n$ :

\begin{itemize}
\item[$\bullet~~$]$R_n$ l'évènement : \og L'athlète réussit à franchir la haie lors de la $n$-ième séance \fg,
\item[$\bullet~~$]$p_n$ la probabilité de l'évènement $R_n$. On considère que $p_0 = 0,6$.
\end{itemize}

\medskip

\begin{enumerate}
\item Soit $n$ un entier naturel, recopier l'arbre pondéré ci-dessous et compléter les
pointillés.

\begin{center}
\pstree[treemode=R,nodesepA=0pt,nodesepB=2.5pt,treesep = 1cm,levelsep=2.5cm]{\TR{}}
{\pstree{\TR{$R_n$~}\taput{$p_n$}}
{\TR{$R_{n+1}$}\taput{\ldots}
\TR{$\overline{R_{n+1}}$}\tbput{\ldots}
}
\pstree{\TR{$\overline{R_n}$~}\tbput{$1 - p_n$}}
{\TR{$R_{n+1}$}\taput{\ldots}
\TR{$\overline{R_{n+1}}$}\tbput{\ldots}
}
}
\end{center}
\item Justifier en vous aidant de l'arbre que, pour tout entier naturel $n$, on a :
\[p_{n+1} = 0,6p_n + 0,3.\]

\item On considère la suite $\left(u_n\right)$ définie, pour tout entier naturel $n$, par $u_n = p_n - 0,75$.

	\begin{enumerate}
		\item Démontrer que la suite $\left(u_n\right)$ est une suite géométrique dont on précisera la raison et le premier terme.
		\item Démontrer que, pour tout entier $n$ naturel $n$ :
		\[p_n = 0,75 - 0,15 \times  0,6^n.\]
		
		\item En déduire que la suite $\left(p_n\right)$ est convergente et déterminer sa limite $\ell$.
		\item Interpréter la valeur de $\ell$ dans le cadre de l'exercice.
	\end{enumerate}
\end{enumerate}

\medskip

\textbf{Partie B}

\medskip

Apres de nombreuses séances d'entraînement, l'entraineur estime maintenant que l'athlète franchit chaque haie avec une probabilité de $0,75$ et ce indépendamment d'avoir franchi ou non les haies précédentes.

On note $X$ la variable aléatoire qui donne le nombre de haies franchies par l'athlète à l'issue d'un $400$~mètres haies qui comporte $10$~haies,

\medskip

\begin{enumerate}
\item Préciser la nature et les paramètres de la loi de probabilité suivie par $X$.
\item Déterminer, à $10^{-3}$ près, la probabilité que l'athlète franchisse les $10$~haies.
\item Calculer $p(X \geqslant 9)$, à $10^{-3}$ près,
\end{enumerate}

\bigskip

\end{document}
